\subsection{Sparse Attention}
\label{subsec:sparse-attn}

The full $N \times N$ attention matrix computed in standard MHA is empirically highly sparse: for typical language modeling tasks, the vast majority of attention weights are near zero, with each token concentrating its attention mass on a small subset of informative positions~\cite{child2019generating}. Sparse attention methods exploit this observation by restricting each query to attend only to a structured or learned subset of key--value positions, reducing the quadratic complexity to $O(N\sqrt{N})$, $O(NW)$, or $O(N \log N)$ depending on the sparsity pattern. The resulting computational savings are particularly significant for long-context modeling, where $N$ ranges from 32K to over 1M tokens in modern LLMs.

The field has evolved through four distinct phases: fixed factorized patterns (2019), combinatorial fixed patterns with theoretical guarantees (2020), hand-crafted heterogeneous layer-wise mixtures deployed in production (2023--2025), and learnable dynamic sparsity that adapts to input content (2024--2025). Throughout this evolution, a recurring theme is the tension between expressivity (the ability to capture arbitrary long-range dependencies) and hardware efficiency (the ability to map sparse access patterns onto contiguous GPU memory operations). We use $N$ for sequence length, $d$ for head dimension, $W$ for window size, $B$ for block size, $k$ for the number of selected blocks, $g$ for the number of global tokens, and $r$ for the number of random connections throughout this subsection.


% ================================================================
\subsubsection{Fixed-Pattern Sparsity}
\label{subsubsec:fixed-pattern}

The earliest approaches to sparse attention replace the dense attention matrix with a predetermined sparsity pattern that is fixed across all inputs and all layers. The central design question is how to choose this pattern so that information can still propagate globally through the composition of multiple layers, despite each individual layer having only local or structured connectivity.

\paragraph{Sparse Transformer.}
Child et al.~\cite{child2019generating} introduced the first systematic study of sparse attention patterns for autoregressive generation. The key idea is to factorize the full attention into two complementary sparse patterns: a \emph{strided} pattern in which each token attends to every $\ell$-th previous token (capturing periodic long-range dependencies), and a \emph{fixed} pattern in which each token attends to a contiguous local window plus a set of fixed ``summary'' positions. Formally, for a sequence of length $N$ and stride $\ell = \lfloor\sqrt{N}\rfloor$, the strided pattern defines the attention set for position $i$ as:
\begin{equation}
\mathcal{A}_{\text{strided}}(i) = \{j : (i - j) \bmod \ell = 0,\; j \leq i\},
\label{eq:sparse-strided}
\end{equation}
yielding $|\mathcal{A}(i)| = O(\sqrt{N})$ per position and $O(N\sqrt{N})$ total complexity. The fixed pattern replaces the stride with summary tokens: each block of $\ell$ consecutive positions designates its last $c$ tokens as aggregators, and all other tokens attend within their block and to these aggregators. By composing strided and fixed patterns across consecutive layers, information can propagate from any position to any other position within two layers, preserving the global receptive field of full attention.

The Sparse Transformer demonstrated that autoregressive models with $O(N\sqrt{N})$
attention can generate high-quality images (CIFAR-10 at 2.80 bits/dim), text (Enwik8 at
0.99 bits/char), and raw audio waveforms with global coherence over sequences of up to
65K time steps.
The accompanying engineering contributions---custom block-sparse CUDA kernels, gradient
checkpointing, and mixed-precision training---were equally influential, establishing the
template for hardware-aware sparse attention implementations that would be followed by
all subsequent work.
The Sparse Transformer's core insight---that sparsity patterns should be designed so that
multi-layer composition recovers global connectivity---remains the guiding principle for
the entire sparse attention research program.

\paragraph{Longformer.}
Beltagy et al.~\cite{beltagy2020longformer} proposed a more practical sparsity pattern tailored to document-level NLP tasks. The central observation is that most tokens require only local context for syntactic processing, while a small number of task-specific tokens (e.g., the \texttt{[CLS]} token for classification, or question tokens for QA) need a global view of the entire sequence. Longformer combines three attention components: (1)~a \emph{sliding window} of size $W$ centered on each token, capturing local context; (2)~a \emph{dilated sliding window} that extends the receptive field by skipping positions at regular intervals with gap $d_g$, expanding the effective window to $W \times d_g$ without additional computation; and (3)~a small number of \emph{global tokens} ($g$ positions) that attend to and are attended by all positions in the sequence. The per-layer complexity is:
\begin{equation}
C_{\text{Longformer}} = O(N \cdot W) + O(N \cdot g),
\label{eq:longformer-complexity}
\end{equation}
which is linear in $N$ when $W$ and $g$ are constants.

A key practical insight is that different layers can use different window sizes---smaller
windows in lower layers (e.g., $W = 256$) and larger windows in upper layers (e.g.,
$W = 512$)---creating a pyramidal receptive field that mirrors the hierarchical nature of
linguistic structure.
Stacking $L$ layers with window size $W$ yields a theoretical receptive field of
$L \times W$ tokens, which for $L = 12$ and $W = 512$ covers 6,144 positions.
The global token mechanism is particularly elegant: by making the global positions
task-dependent, Longformer bridges the gap between generic sparse attention and
task-specific information routing.
On long-document benchmarks (WikiHop, TriviaQA, HotpotQA), Longformer consistently
outperformed RoBERTa with truncated inputs, demonstrating the practical value of
processing full documents rather than truncating to 512 tokens.
The sliding window plus global token design became the conceptual template for subsequent
production deployments: Mistral adopted the sliding window, Gemma adopted the
local/global layer distinction, and both BigBird and later models retained the idea that
a small number of globally-connected positions can compensate for the limited range of
local windows.

\paragraph{BigBird.}
Zaheer et al.~\cite{zaheer2020big} extended the fixed-pattern approach by adding a \emph{random} attention component to the sliding window and global token mechanisms. Each token additionally attends to $r$ randomly selected positions, yielding the combined attention pattern:
\begin{equation}
\mathcal{A}_{\text{BigBird}}(i) = \underbrace{\mathcal{A}_{\text{window}}(i)}_{\text{local}} \;\cup\; \underbrace{\mathcal{A}_{\text{random}}(i)}_{\text{random}} \;\cup\; \underbrace{\mathcal{A}_{\text{global}}}_{\text{global}},
\label{eq:bigbird}
\end{equation}
where $\mathcal{A}_{\text{random}}(i)$ contains $r$ randomly sampled positions, $\mathcal{A}_{\text{window}}(i)$ is a local window of size $W$, and $\mathcal{A}_{\text{global}}$ is a fixed set of $g$ global positions. The random component is motivated by the theory of random graphs: Erd\H{o}s--R\'{e}nyi graphs with $O(\log N)$ random edges per node have $O(\log N)$ diameter, ensuring that information can propagate between any two tokens in logarithmically many layers.

BigBird's most significant contribution is theoretical: the authors proved that the BigBird attention pattern is a \emph{universal approximator} of continuous sequence-to-sequence functions and is Turing complete, establishing that $O(N)$ attention connections per layer suffice for universal computation provided the sparsity pattern includes both local and global components. This result provides a formal justification for the entire sparse attention research program: the quadratic cost of full attention is not a necessary price for expressiveness, but rather an artifact of the specific dense connectivity pattern. The overall complexity is $O(N \cdot (W + r + g))$, linear in $N$ for fixed hyperparameters. In practice, BigBird achieved state-of-the-art results on long-document tasks with sequences up to 4,096 tokens.

The fixed-pattern methods share a common limitation: the sparsity structure is determined at design time and cannot adapt to the input. A sliding window of size 512 treats all tokens identically, regardless of whether the current query requires local syntactic information or long-range coreference resolution. This rigidity motivated the development of data-dependent sparsity methods.


% ================================================================
\subsubsection{Hash-Based Sparsity}
\label{subsubsec:hash-sparsity}

Fixed-pattern methods impose the same sparsity structure regardless of input content, which is suboptimal when the distribution of informative token pairs varies across examples. Data-dependent sparsity methods address this by dynamically identifying which query--key pairs are likely to have high attention weights.

\paragraph{Reformer (LSH Attention).}
Kitaev et al.~\cite{kitaev2020reformer} took a fundamentally different approach by using \emph{locality-sensitive hashing} (LSH) to identify, for each query, the keys that are most likely to receive high attention weights. The core idea is that if $q_i$ and $k_j$ have a high dot product, they are likely to hash to the same bucket under a random hyperplane projection. The hash function is defined as:
\begin{equation}
h(x) = \arg\max\bigl([xR;\; -xR]\bigr), \quad R \in \mathbb{R}^{d \times b/2},
\label{eq:lsh-hash}
\end{equation}
where $R$ is a random projection matrix and $b$ is the number of buckets. After hashing, tokens are sorted by bucket, and attention is computed only within each bucket (with an additional local window to handle bucket boundaries). Using $n_r$ rounds of independent hashing ($n_r = 8$ is recommended) to reduce the probability of missing high-attention pairs, the expected complexity is $O(N \cdot n_r \cdot N/b)$, which reduces to $O(N \log N)$ when the bucket size is set to $b = O(N / \log N)$.

A subtle but important design choice in Reformer is the use of \emph{shared query--key
projections} ($Q = K$), which ensures that $q_i \cdot k_j = q_i \cdot q_j$ and thus
tokens that are similar in query space are also similar in key space.
This sharing is necessary for the LSH approximation to work well: if $Q$ and $K$ were
independent projections, high dot-product pairs in $QK^\top$ would not necessarily
correspond to nearby points in the hash space.
The cost of this constraint is a slight reduction in the model's representational
flexibility, as the query and key projections can no longer learn independent
transformations of the input.

Reformer also introduced reversible residual layers (RevNet), which eliminate the need
to store intermediate activations during backpropagation by computing
$y_1 = x_1 + F(x_2)$, $y_2 = x_2 + G(y_1)$, reducing activation memory from
$O(N \cdot L)$ to $O(N)$.
This dual contribution---algorithmic sparsity plus memory-efficient
backpropagation---made Reformer the first model to demonstrate training on sequences
of 64K tokens, a significant milestone at the time.

Despite its theoretical elegance, LSH attention has several practical limitations that explain its limited adoption in production LLMs. The multi-round hashing introduces significant constant-factor overhead, effectively multiplying the per-token cost by $n_r$. The sorting operation disrupts the regular memory access patterns that GPUs are optimized for, resulting in wall-clock speedups that fall far short of the theoretical $O(N \log N / N^2)$ ratio, particularly for the moderate sequence lengths ($N \leq 8{,}192$) common in practice. The approximation quality degrades when the attention distribution is not sufficiently peaked, which is common in lower layers of deep Transformers where attention tends to be more diffuse. These factors led the community to gravitate toward simpler fixed-pattern methods for production use, though Reformer's conceptual contribution---that data-dependent sparsity can be achieved through hashing---directly influenced subsequent work on learnable routing mechanisms such as NSA and MoBA.


% ================================================================
\subsubsection{Heterogeneous Layer-Wise Mixtures in Production}
\label{subsubsec:heterogeneous-sparse}

The second generation of sparse attention, deployed in production-scale LLMs from 2023 onward, abandons the goal of a single universal sparsity pattern in favor of \emph{heterogeneous} configurations that assign different attention types to different layers. This approach is motivated by the empirical observation that shallow layers primarily capture local syntactic patterns (for which a small window suffices), while deeper layers require broader context for semantic reasoning. The key engineering insight is that simple patterns---particularly the sliding window---are far easier to implement efficiently on GPU hardware than the more sophisticated patterns proposed in the academic literature.

\paragraph{Sliding Window Attention (SWA).}
Mistral 7B~\cite{jiang2023mistral} demonstrated that a simple sliding window of size $W = 4{,}096$ applied uniformly across all layers can achieve competitive quality with full attention for models up to 7B parameters. The key insight is that information propagates across windows through the stacking of layers: after $L$ layers, the effective receptive field is $L \times W$ tokens, which for $L = 32$ and $W = 4{,}096$ covers 131K tokens---far exceeding the training context length. SWA reduces the per-layer complexity from $O(N^2)$ to $O(NW)$ and, critically, caps the KV-cache size at $W$ entries per layer regardless of the total sequence length. The \emph{rolling buffer KV cache} implements this as a fixed-size circular buffer where position $i$ is stored at index $i \bmod W$, enabling constant-memory streaming inference over arbitrarily long inputs. Combined with grouped-query attention (GQA with 8 KV heads for 32 query heads),
Mistral~7B established SWA as the most hardware-friendly sparse attention variant,
requiring no custom kernels beyond a simple causal mask modification.

It is worth noting that SWA provides different benefits during the two phases of
autoregressive inference.
During \emph{prefill} (processing the prompt), SWA reduces the quadratic attention
computation from $O(N^2)$ to $O(NW)$, which is significant for long prompts.
During \emph{decode} (generating tokens one at a time), the primary benefit is memory
rather than compute: the rolling buffer caps the KV-cache memory at $O(W \cdot L \cdot d)$
regardless of the total generated length, preventing the linear memory growth that
otherwise limits the maximum sequence length during serving.
This memory-bounded property is particularly valuable for streaming applications and
multi-turn conversations where the total context can grow unboundedly.

\paragraph{Interleaved Local/Global Attention.}
Gemma~2~\cite{gemma2024improving} introduced a more nuanced approach by interleaving sliding-window (local) and full (global) attention layers in a 1:1 ratio. The local layers handle fine-grained syntactic processing with $O(NW)$ cost, while the global layers provide the cross-sequence information flow needed for long-range reasoning at $O(N^2)$ cost. Since only half the layers compute full attention, the amortized complexity is $O(N^2/2 + NW/2)$---a meaningful reduction for long sequences.

Gemma~3~\cite{gemma2025gemma3} sharpened this ratio to 5:1 (local-to-global) with a reduced local window ($W = 1{,}024$) and, crucially, distinct RoPE base frequencies for each layer type: $\theta_{\text{local}} = 10{,}000$ for local layers and $\theta_{\text{global}} = 1{,}000{,}000$ for global layers. This frequency differentiation allows global layers to encode much longer-range positional information without degrading the local layers' resolution for nearby tokens. The 5:1 ratio means that only ${\sim}17\%$ of layers incur quadratic cost, making the architecture substantially more efficient for long contexts while maintaining quality within 1\% of full-attention baselines. The independent convergence of Gemma~3's 5:1 ratio with the attention-to-recurrence ratios observed in hybrid architectures (Section~\ref{subsec:hybrid}) suggests a natural balance point between local and global processing that may reflect the fraction of layers requiring precise long-range retrieval.

\paragraph{Temporal Convolution on KV.}
Baichuan-M1~\cite{baichuan2025m1} introduced a complementary technique: applying a short causal depthwise one-dimensional convolution (kernel size 3--7) over the key and value representations before attention computation. The convolved representations $K' = \text{CausalConv1D}(K)$ and $V' = \text{CausalConv1D}(V)$ fuse information from adjacent tokens along the sequence dimension, effectively injecting local context into the KV representations themselves rather than relying solely on the attention mask. The causality constraint ensures that position $t$ depends only on positions $\leq t$, maintaining compatibility with autoregressive generation. This adds less than 1\% additional FLOPs while improving the model's ability to capture fine-grained temporal patterns---particularly valuable in the medical domain targeted by Baichuan-M1, where drug names, diagnostic codes, and numerical sequences require precise local pattern matching. The design echoes the short convolution branches in Mamba~\cite{gu2024mamba} and H3, but embeds them \emph{within} the attention mechanism rather than as a replacement, preserving full compatibility with standard Transformer infrastructure and KV-cache-based inference.

\paragraph{Design considerations.}
The choice of window size $W$ involves a three-way trade-off among quality, memory, and latency. Empirically, $W = 4{,}096$ (Mistral) provides strong general-purpose performance, while $W = 1{,}024$ (Gemma~3) is viable when interleaved with global layers. The local-to-global layer ratio determines how much of the model's capacity is devoted to long-range reasoning versus local processing; the convergence of multiple independent efforts on ratios between 1:1 and 5:1 suggests a robust operating regime. An emerging best practice is to use different positional encoding configurations for local and global layers, as the optimal frequency spectrum for short-range and long-range position encoding differs substantially. The heterogeneous approach has become the dominant strategy in production LLMs as of 2025. Table~\ref{tab:sparse-production} summarizes the configurations adopted by major models.

\begin{table}[t]
\centering
\caption{Sparse attention configurations in production LLMs. $W$: sliding window size. The local:global ratio indicates the fraction of layers using each attention type.}
\label{tab:sparse-production}
\small
\begin{tabular}{lcccl}
\toprule
\textbf{Model} & \textbf{Local:Global} & \textbf{Window $W$} & \textbf{RoPE Diff.} & \textbf{Additional} \\
\midrule
Mistral 7B     & All local    & 4,096  & No  & Rolling buffer KV cache \\
Gemma 2        & 1:1          & 4,096  & No  & Logit soft-capping \\
Gemma 3        & 5:1          & 1,024  & Yes ($\theta$: 10K/1M) & --- \\
Yi-Lightning   & Hybrid       & Varies & No  & Cross-layer KV sharing \\
Baichuan-M1    & Hybrid       & Varies & No  & Temporal conv on KV \\
\bottomrule
\end{tabular}
\end{table}


% ================================================================
\subsubsection{Learnable Dynamic Sparsity}
\label{subsubsec:learnable-sparse}

The most recent phase of sparse attention development (2024--2025) replaces hand-crafted patterns with \emph{learned} sparsity that adapts dynamically to the input content. The key motivation is that the optimal set of tokens to attend to varies not only across layers but across inputs and even across individual query positions within a single input. This shift parallels the evolution from fixed expert assignment to learned routing in Mixture-of-Experts architectures and represents a fundamental change in how attention budgets are allocated. Three recent approaches exemplify this direction, each drawing on a different mechanism for learning which tokens to attend to.

\paragraph{Native Sparse Attention (NSA).}
Yuan et al.~\cite{yuan2025native}, developed at DeepSeek, introduced a hardware-aligned, natively trainable three-branch architecture. NSA decomposes attention into three complementary branches---compression, selection, and sliding window---and aggregates their outputs via learned gates:
\begin{equation}
o_i = g_1^{(i)} \cdot o_{\text{compress}}^{(i)} + g_2^{(i)} \cdot o_{\text{select}}^{(i)} + g_3^{(i)} \cdot o_{\text{window}}^{(i)},
\label{eq:nsa}
\end{equation}
where the gates $g_1, g_2, g_3$ are per-head scalars produced by a linear projection followed by sigmoid activation, modulated by the input. The \emph{compression branch} partitions the KV sequence into blocks of size $B$ and compresses each block into a single summary token via mean pooling, providing a coarse global view at $O(N/B)$ cost. The \emph{selection branch} uses the compressed attention scores to identify the top-$k$ most relevant blocks, then computes fine-grained attention only within those blocks, at $O(kB)$ cost per query. The \emph{sliding window branch} attends to the most recent $W$ tokens for local context. The total per-query complexity is:
\begin{equation}
C_{\text{NSA}} = O\!\left(\frac{N}{B} + kB + W\right),
\label{eq:nsa-complexity}
\end{equation}
which is linear in $N$ when $k$ and $B$ are treated as constants.

A critical design principle is \emph{hardware alignment}: all three branches operate on contiguous memory blocks, ensuring that the sparse access patterns map efficiently onto GPU tensor cores. The block selection mechanism is implemented via Triton kernels that maintain coalesced memory access patterns, directly addressing the hardware inefficiency that plagued Reformer's LSH-based approach. NSA is trained end-to-end with the sparsity pattern emerging from the learned gates, rather than being imposed as a fixed prior. Empirically, NSA achieves quality comparable to full attention on language modeling benchmarks while reducing the effective attention computation by 4--8$\times$ for sequences of 64K tokens, with the compression branch capturing global trends, the selection branch focusing on salient long-range dependencies, and the window branch handling local syntax.

\paragraph{Mixture of Block Attention (MoBA).}
Lu et al.~\cite{lu2025moba}, developed at Moonshot AI (Kimi), drew a direct analogy between sparse attention and Mixture-of-Experts: just as MoE routes each token to the most relevant expert FFN layers, MoBA routes each query to the most relevant KV blocks. The KV sequence is partitioned into $B_{\text{total}}$ blocks, and each query token routes to the top-$k$ most relevant blocks via a lightweight router that computes relevance scores as the dot product between the query and each block's mean key:
\begin{equation}
s_{i,j} = q_i \cdot \bar{k}_j, \quad \bar{k}_j = \frac{1}{|\mathcal{B}_j|}\sum_{t \in \mathcal{B}_j} k_t,
\label{eq:moba-router}
\end{equation}
where $\bar{k}_j$ is the mean key vector of block $j$. Each query selects the top-$k$ blocks with the highest scores and computes standard attention within them, supplemented by a causal sliding window to ensure local coverage. MoBA achieves sub-quadratic complexity ($O(NkB)$ with $kB \ll N$) while preserving the mathematical equivalence to full attention when $k$ equals the total number of blocks.

A notable property is that MoBA is a strict generalization of sliding window attention: if the router always selects the blocks closest to the query position, MoBA degenerates to SWA. In practice, MoBA always includes the most recent block (analogous to a sliding window) and routes to $k - 1$ additional blocks, combining the reliability of local attention with the flexibility of content-based routing. On 128K-context benchmarks, MoBA achieves performance close to full attention with substantially reduced computation. A key practical advantage is seamless compatibility with existing full-attention models: during training, some layers can use MoBA while others retain full attention, and the routing mechanism can be toggled off to recover exact full attention for verification or fine-tuning.

\paragraph{Mixture of Attention (MoA).}
Fu et al.~\cite{fu2024moa} addressed a complementary question: rather than designing a single sparsity pattern for all layers and heads, can we automatically assign the optimal pattern to each head? MoA constructs a search space of candidate attention patterns---full attention, sliding window (with varying sizes), local-plus-global, block sparse, and static sparse---and uses profiling on calibration data to measure the approximation error of each pattern for each head. An evolutionary search under a given compute budget then selects the per-head configuration that minimizes total approximation error.

Applied to Llama-2 models (7B and 13B), MoA reduces attention computation by 50--70\% with less than 1\% quality degradation on LongBench, and reveals that the optimal pattern varies substantially across layers: lower layers tend to require more full attention heads, while upper layers can tolerate aggressive sparsity. Unlike NSA and MoBA, which learn sparsity patterns during pretraining, MoA operates as a \emph{post-training} compression method, making it applicable to any pretrained model without retraining. With compatible sparse attention kernels, MoA achieves 1.5--2$\times$ inference speedup on long-context workloads, validating the hypothesis that heterogeneous per-head sparsity is strictly superior to uniform sparsity. This finding provides empirical support for the heterogeneous layer designs adopted by Gemma~2/3 and suggests that automated sparsity search may become a standard post-training optimization step.

\paragraph{NSA vs.\ MoBA: complementary design philosophies.}
Although NSA and MoBA both achieve learnable dynamic sparsity, they embody different
design philosophies.
NSA uses a \emph{multi-resolution} approach: the compression branch provides a coarse
global summary, the selection branch refines it at block granularity, and the window
branch supplies local detail.
The three branches operate in parallel and are fused through learned gates, giving each
head the flexibility to weight global, selective, and local information differently.
MoBA, by contrast, uses a \emph{flat routing} approach: all blocks are treated as
equal-sized candidates, and the router selects the top-$k$ most relevant ones.
This design is simpler and more directly analogous to MoE, but does not provide the
multi-resolution hierarchy of NSA.
In practice, the two approaches achieve comparable quality--efficiency trade-offs,
suggesting that the specific mechanism for selecting which tokens to attend to matters
less than the principle of making the selection learned and hardware-aligned.
Both methods share the critical property of block-contiguous memory access, and both
retain a sliding window for local context.
The emergence of these two approaches from independent research groups (DeepSeek and
Moonshot AI) suggests that the block-based, hardware-aligned paradigm for learnable
sparse attention is robust and likely to become the standard framework for future work.


% ================================================================
\subsubsection{Comparison and Analysis}
\label{subsubsec:sparse-comparison}

Table~\ref{tab:sparse-attn-comparison} provides a comprehensive comparison of sparse attention methods across the four phases of development.

\begin{table}[t]
\centering
\caption{Comparison of sparse attention methods. $N$: sequence length, $W$: window size, $B$: block size, $k$: number of selected blocks, $g$: number of global tokens, $r$: number of random connections, $R$: local-to-global layer ratio.}
\label{tab:sparse-attn-comparison}
\small
\begin{tabular}{lccccl}
\toprule
\textbf{Method} & \textbf{Year} & \textbf{Complexity} & \textbf{Pattern Type} & \textbf{Trainable} & \textbf{Representative Models} \\
\midrule
Sparse Transformer & 2019 & $O(N\sqrt{N})$ & Fixed (strided) & No & --- \\
Longformer         & 2020 & $O(NW + Ng)$   & Fixed (window+global) & No & --- \\
BigBird            & 2020 & $O(NW + Ng + Nr)$ & Fixed (window+global+random) & No & --- \\
Reformer           & 2020 & $O(N \log N)$  & Hash-based (LSH) & No & --- \\
Mistral SWA        & 2023 & $O(NW)$        & Fixed (window) & No & Mistral 7B \\
Gemma 2/3          & 2024 & $O(NW + N^2/R)$ & Heterogeneous (interleaved) & No & Gemma 2, Gemma 3 \\
NSA                & 2025 & $O(N/B{+}kB{+}W)$ & Learned (3-branch gated) & Yes & DeepSeek \\
MoBA               & 2025 & $O(NkB)$       & Learned (block routing) & Yes & Kimi \\
MoA                & 2024 & varies         & Searched (per-head) & Post-hoc & Llama-2 (compressed) \\
\bottomrule
\end{tabular}
\end{table}

The evolution from fixed-pattern to learnable sparsity reveals several recurring themes that provide guidance for future sparse attention design.

\emph{Locality is the strongest prior.} Every successful sparse attention method retains a sliding window or local attention component, reflecting the empirical observation that adjacent tokens are almost always relevant. The methods differ primarily in how they handle long-range dependencies on top of this local base---through global tokens (Longformer, BigBird), hash-based grouping (Reformer), interleaved global layers (Gemma), or learned block selection (NSA, MoBA). This universality of the local window suggests that it captures a fundamental property of natural language: the strong locality of syntactic and short-range semantic dependencies. Any future sparse attention design should treat local attention as a non-negotiable baseline.

\emph{Hardware alignment is as important as algorithmic complexity.} The practical speedup of a sparse attention method depends not only on the number of non-zero entries but on the \emph{contiguity} of memory accesses. Reformer's $O(N \log N)$ theoretical complexity did not translate into practical speedups because its irregular memory access patterns---sorting, bucketing, and gathering---are poorly suited to GPU architectures optimized for coalesced, sequential memory access. In contrast, Mistral's simpler $O(NW)$ sliding window achieves substantial wall-clock savings through regular, predictable memory access patterns that map directly to FlashAttention-style tiling~\cite{dao2022flashattention}. NSA explicitly addresses this lesson by designing all operations around contiguous block access patterns, and MoBA's block-level routing similarly maintains memory locality. This principle---that sparse attention patterns must be co-designed with the memory hierarchy of the target hardware---is likely to become even more important as models scale to longer contexts and are deployed on diverse accelerator architectures.

\emph{Heterogeneous configurations outperform uniform ones.} The transition from homogeneous patterns (all layers use the same window) to heterogeneous configurations---pioneered by Gemma~2/3's interleaved local/global layers and validated by MoA's automated per-head search---reveals that the optimal attention pattern varies across the depth of the network. Lower layers, which build local syntactic representations, benefit from dense local attention, while upper layers, which perform more abstract reasoning, can tolerate sparser patterns or benefit from selective global attention. This observation connects sparse attention to the broader theme of layer-wise specialization that also appears in hybrid architectures (Section~\ref{subsec:hybrid}) and suggests that future models will increasingly adopt per-layer or even per-head attention configurations.

\emph{From fixed to learned sparsity.} The convergence of MoE-style routing (MoBA), multi-branch gating (NSA), and automated search (MoA) points toward a future where sparse attention patterns are fully learned end-to-end, with each layer and each head independently determining its attention budget based on the input. This trajectory mirrors the broader deep learning trend of replacing hand-crafted inductive biases with learned ones, and suggests that the distinction between ``sparse attention'' and ``full attention'' may eventually dissolve into a continuum of learned attention allocation strategies.

The relationship between sparse attention and other token mixer families is complementary rather than competitive. Sparse attention composes naturally with KV-cache optimization (e.g., GQA + sliding window in Gemma~3), with KV-cache compression (e.g., evicting tokens that receive zero attention weight), and with linear attention or SSM layers in hybrid architectures (Section~\ref{subsec:hybrid}). The combination of sparse attention with temporal convolution on KV (Baichuan-M1) further illustrates that sparsity in the attention mask and local feature extraction in the representation space are orthogonal and complementary mechanisms. Open challenges include: (i)~developing efficient training algorithms that can learn sparsity patterns without requiring full attention as a teacher signal; (ii)~extending dynamic sparsity to the prefill phase, where the quadratic cost is most acute; (iii)~establishing theoretical frameworks for selecting optimal sparsity configurations---window sizes, local-to-global ratios, and block sizes---that go beyond empirical ablation; and (iv)~co-designing sparse attention kernels with emerging hardware architectures to close the gap between theoretical and realized throughput.
